% Options for packages loaded elsewhere
\PassOptionsToPackage{unicode}{hyperref}
\PassOptionsToPackage{hyphens}{url}
\documentclass[
]{article}
\usepackage{xcolor}
\usepackage{amsmath,amssymb}
\setcounter{secnumdepth}{-\maxdimen} % remove section numbering
\usepackage{iftex}
\ifPDFTeX
  \usepackage[T1]{fontenc}
  \usepackage[utf8]{inputenc}
  \usepackage{textcomp} % provide euro and other symbols
\else % if luatex or xetex
  \usepackage{unicode-math} % this also loads fontspec
  \defaultfontfeatures{Scale=MatchLowercase}
  \defaultfontfeatures[\rmfamily]{Ligatures=TeX,Scale=1}
\fi
\usepackage{lmodern}
\ifPDFTeX\else
  % xetex/luatex font selection
\fi
% Use upquote if available, for straight quotes in verbatim environments
\IfFileExists{upquote.sty}{\usepackage{upquote}}{}
\IfFileExists{microtype.sty}{% use microtype if available
  \usepackage[]{microtype}
  \UseMicrotypeSet[protrusion]{basicmath} % disable protrusion for tt fonts
}{}
\makeatletter
\@ifundefined{KOMAClassName}{% if non-KOMA class
  \IfFileExists{parskip.sty}{%
    \usepackage{parskip}
  }{% else
    \setlength{\parindent}{0pt}
    \setlength{\parskip}{6pt plus 2pt minus 1pt}}
}{% if KOMA class
  \KOMAoptions{parskip=half}}
\makeatother
\usepackage{color}
\usepackage{fancyvrb}
\newcommand{\VerbBar}{|}
\newcommand{\VERB}{\Verb[commandchars=\\\{\}]}
\DefineVerbatimEnvironment{Highlighting}{Verbatim}{commandchars=\\\{\}}
% Add ',fontsize=\small' for more characters per line
\newenvironment{Shaded}{}{}
\newcommand{\AlertTok}[1]{\textcolor[rgb]{1.00,0.00,0.00}{\textbf{#1}}}
\newcommand{\AnnotationTok}[1]{\textcolor[rgb]{0.38,0.63,0.69}{\textbf{\textit{#1}}}}
\newcommand{\AttributeTok}[1]{\textcolor[rgb]{0.49,0.56,0.16}{#1}}
\newcommand{\BaseNTok}[1]{\textcolor[rgb]{0.25,0.63,0.44}{#1}}
\newcommand{\BuiltInTok}[1]{\textcolor[rgb]{0.00,0.50,0.00}{#1}}
\newcommand{\CharTok}[1]{\textcolor[rgb]{0.25,0.44,0.63}{#1}}
\newcommand{\CommentTok}[1]{\textcolor[rgb]{0.38,0.63,0.69}{\textit{#1}}}
\newcommand{\CommentVarTok}[1]{\textcolor[rgb]{0.38,0.63,0.69}{\textbf{\textit{#1}}}}
\newcommand{\ConstantTok}[1]{\textcolor[rgb]{0.53,0.00,0.00}{#1}}
\newcommand{\ControlFlowTok}[1]{\textcolor[rgb]{0.00,0.44,0.13}{\textbf{#1}}}
\newcommand{\DataTypeTok}[1]{\textcolor[rgb]{0.56,0.13,0.00}{#1}}
\newcommand{\DecValTok}[1]{\textcolor[rgb]{0.25,0.63,0.44}{#1}}
\newcommand{\DocumentationTok}[1]{\textcolor[rgb]{0.73,0.13,0.13}{\textit{#1}}}
\newcommand{\ErrorTok}[1]{\textcolor[rgb]{1.00,0.00,0.00}{\textbf{#1}}}
\newcommand{\ExtensionTok}[1]{#1}
\newcommand{\FloatTok}[1]{\textcolor[rgb]{0.25,0.63,0.44}{#1}}
\newcommand{\FunctionTok}[1]{\textcolor[rgb]{0.02,0.16,0.49}{#1}}
\newcommand{\ImportTok}[1]{\textcolor[rgb]{0.00,0.50,0.00}{\textbf{#1}}}
\newcommand{\InformationTok}[1]{\textcolor[rgb]{0.38,0.63,0.69}{\textbf{\textit{#1}}}}
\newcommand{\KeywordTok}[1]{\textcolor[rgb]{0.00,0.44,0.13}{\textbf{#1}}}
\newcommand{\NormalTok}[1]{#1}
\newcommand{\OperatorTok}[1]{\textcolor[rgb]{0.40,0.40,0.40}{#1}}
\newcommand{\OtherTok}[1]{\textcolor[rgb]{0.00,0.44,0.13}{#1}}
\newcommand{\PreprocessorTok}[1]{\textcolor[rgb]{0.74,0.48,0.00}{#1}}
\newcommand{\RegionMarkerTok}[1]{#1}
\newcommand{\SpecialCharTok}[1]{\textcolor[rgb]{0.25,0.44,0.63}{#1}}
\newcommand{\SpecialStringTok}[1]{\textcolor[rgb]{0.73,0.40,0.53}{#1}}
\newcommand{\StringTok}[1]{\textcolor[rgb]{0.25,0.44,0.63}{#1}}
\newcommand{\VariableTok}[1]{\textcolor[rgb]{0.10,0.09,0.49}{#1}}
\newcommand{\VerbatimStringTok}[1]{\textcolor[rgb]{0.25,0.44,0.63}{#1}}
\newcommand{\WarningTok}[1]{\textcolor[rgb]{0.38,0.63,0.69}{\textbf{\textit{#1}}}}
\usepackage{longtable,booktabs,array}
\newcounter{none} % for unnumbered tables
\usepackage{calc} % for calculating minipage widths
% Correct order of tables after \paragraph or \subparagraph
\usepackage{etoolbox}
\makeatletter
\patchcmd\longtable{\par}{\if@noskipsec\mbox{}\fi\par}{}{}
\makeatother
% Allow footnotes in longtable head/foot
\IfFileExists{footnotehyper.sty}{\usepackage{footnotehyper}}{\usepackage{footnote}}
\makesavenoteenv{longtable}
\usepackage{graphicx}
\makeatletter
\newsavebox\pandoc@box
\newcommand*\pandocbounded[1]{% scales image to fit in text height/width
  \sbox\pandoc@box{#1}%
  \Gscale@div\@tempa{\textheight}{\dimexpr\ht\pandoc@box+\dp\pandoc@box\relax}%
  \Gscale@div\@tempb{\linewidth}{\wd\pandoc@box}%
  \ifdim\@tempb\p@<\@tempa\p@\let\@tempa\@tempb\fi% select the smaller of both
  \ifdim\@tempa\p@<\p@\scalebox{\@tempa}{\usebox\pandoc@box}%
  \else\usebox{\pandoc@box}%
  \fi%
}
% Set default figure placement to htbp
\def\fps@figure{htbp}
\makeatother
\setlength{\emergencystretch}{3em} % prevent overfull lines
\providecommand{\tightlist}{%
  \setlength{\itemsep}{0pt}\setlength{\parskip}{0pt}}
\usepackage{bookmark}
\IfFileExists{xurl.sty}{\usepackage{xurl}}{} % add URL line breaks if available
\urlstyle{same}
\hypersetup{
  pdftitle={02\_04\_Tag\_2\_Mailcow\_AD\_Integration},
  hidelinks,
  pdfcreator={LaTeX via pandoc}}

\title{02\_04\_Tag\_2\_Mailcow\_AD\_Integration}
\author{}
\date{}

\begin{document}
\maketitle

\subsection{Windows-Linux Integration
(Mailserver)}\label{windows-linux-integration-mailserver}

\textbf{Datum:} 10.03.2026\\
\textbf{Bearbeiter:} Samuel

\subsection{👥 Gruppenarbeit}\label{gruppenarbeit}

\begin{enumerate}
\tightlist
\item
  \textbf{Infrastruktur-Abstimmung:} Festlegung des Service-Accounts für
  den LDAP-Bind im Windows-Team.
\item
  \textbf{DNS-Troubleshooting:} Behebung eines Auflösungsfehlers
  (beta.local wurde intern nicht korrekt auf den DC geroutet).
\end{enumerate}

\subsection{2. Technische Umsetzung
(Detail-Dokumentation)}\label{technische-umsetzung-detail-dokumentation}

\subsubsection{2.1 Unbound
DNS-Konfiguration}\label{unbound-dns-konfiguration}

Damit der Mailserver die Domäne \texttt{net13.beta} und den DC
\texttt{192.168.13.10} finden kann, musste der interne Resolver
angepasst werden.

\textbf{Datei:} \texttt{data/conf/unbound/unbound.conf}

\begin{Shaded}
\begin{Highlighting}[]
\NormalTok{server:}
\NormalTok{  \# ... (bestehende Config)}
  
\NormalTok{  \# Stub{-}Zone für lokale Domäne}
\NormalTok{  local{-}zone: "net13.beta." transparent}
\NormalTok{  domain{-}insecure: "net13.beta"}

\NormalTok{stub{-}zone:}
\NormalTok{  name: "net13.beta"}
\NormalTok{  stub{-}addr: 192.168.13.10  \# IP des Domain Controllers}
\end{Highlighting}
\end{Shaded}

\emph{Aktion:} \texttt{docker\ compose\ restart\ unbound-mailcow}

\subsubsection{2.2 LDAP-Anbindung (Mailcow
UI)}\label{ldap-anbindung-mailcow-ui}

Die Authentifizierung wurde im Admin-Panel unter
\texttt{System\ \textgreater{}\ Configuration\ \textgreater{}\ Access\ \textgreater{}\ LDAP}
konfiguriert.

{\def\LTcaptype{none} % do not increment counter
\begin{longtable}[]{@{}lll@{}}
\toprule\noalign{}
Parameter & Wert & Erklärung \\
\midrule\noalign{}
\endhead
\bottomrule\noalign{}
\endlastfoot
\textbf{Server Address} & \texttt{192.168.13.10} & IP des DC \\
\textbf{Port} & \texttt{389} (StartTLS) oder \texttt{636} (LDAPS) &
Hier: 636 mit "Skip Verify" (Labor) \\
\textbf{Base DN} & \texttt{DC=net13,DC=beta} & Suchbereich \\
\textbf{Bind DN} & \texttt{CN=svc\_mailcow,OU=IT,DC=net13,DC=beta} &
Service-User zum Lesen \\
\textbf{User Filter} &
\texttt{(\&(objectClass=user)(sAMAccountName=\%u))} & Login via
sAMAccountName \\
\textbf{Sync Interval} & \texttt{300} & Sync alle 5 Min \\
\end{longtable}
}

\subsubsection{2.3 Validierung}\label{validierung}

\begin{enumerate}
\tightlist
\item
  \textbf{Verbindungstest:} Button "Test Connection" -\textgreater{}
  \textbf{OK}.
\item
  \textbf{SOGo Login:} User \texttt{d.zimmermann} (HR) konnte sich mit
  AD-Passwort anmelden.
\item
  \textbf{Log-Check:}
  \texttt{docker\ compose\ logs\ -f\ dovecot-mailcow} zeigte
  erfolgreichen Auth-Handshake.
\end{enumerate}

\subsection{📸 Dokumentation \&
Screenshots}\label{dokumentation-screenshots}

\pandocbounded{\includegraphics[keepaspectratio,alt={Mailcow Container-Status}]{C:/Users/User/Documents/SamuelDenkwerkstätte/_assets/mailcow_container_status.png}}

\begin{itemize}
\tightlist
\item
  \textbf{Screenshot 1:} Status der Mailcow-Container in Docker
  (\texttt{docker\ ps}).\\
  Siehe Abbildung oben.
\item
  \textbf{Screenshot 2:} LDAP-Konfigurationsmaske (Test Connection
  erfolgreich).\\
  \pandocbounded{\includegraphics[keepaspectratio,alt={Mailcow LDAP-Konfigurationsmaske}]{C:/Users/User/Documents/SamuelDenkwerkstätte/_assets/mailcow_ldap_configuration.png}}
\end{itemize}

\subsubsection{Ergänzende Nachweise für
LDAPS}\label{erguxe4nzende-nachweise-fuxfcr-ldaps}

\pandocbounded{\includegraphics[keepaspectratio,alt={Windows Firewall Regelkonfiguration fuer LDAPS / Domain Services}]{C:/Users/User/Documents/SamuelDenkwerkstätte/_assets/windows_firewall_ldaps_rule.png}}

\pandocbounded{\includegraphics[keepaspectratio,alt={LDAPS-Verbindungstest auf Port 636}]{C:/Users/User/Documents/SamuelDenkwerkstätte/_assets/ldp_ldaps_port_636_test.png}}

\subsection{📝 Reflexion / Offene Punkte}\label{reflexion-offene-punkte}

\begin{itemize}
\tightlist
\item
  Die Integration vereinfacht die Benutzerverwaltung erheblich (Single
  Source of Truth im AD).
\item
  \textbf{Lessons Learned:} Bei LDAP über Port 636 muss das Zertifikat
  des DCs im Mailcow-Container bekannt sein (oder "Ignore SSL Errors"
  temporär aktiv).
\item
  Nächster Schritt: Einrichtung der Backup-Strategie für das gesamte
  Mail-System.
\end{itemize}

\end{document}
